\chapter{Over The Counter Medicines}

\medskip
\noindent\textit{\textbf{Under review.} This chapter is an AI-assisted educational draft; it is not a substitute for professional medical advice.}

\section*{Definition and Overview}
Over The Counter Medicines is a medicine or treatment-related topic. Its safe use depends on the indication, dose, interactions, contraindications, and monitoring required for the individual.

\section*{Clinical Features}
Possible effects and adverse effects vary with the medicine and may include gastrointestinal, neurological, allergic, cardiovascular, or organ-specific problems.

\section*{When to Seek Medical Care}
Seek medical review for new, persistent, worsening, or function-limiting symptoms. Call 000 for severe breathing difficulty, collapse, sudden neurological change, severe chest or abdominal pain, or any rapidly worsening illness.

\section*{Diagnosis and Investigations}
Review the indication, medicines, allergies, kidney and liver function, pregnancy status, and relevant monitoring tests.

\section*{Management}
Use only as directed by a qualified clinician or pharmacist. Do not share medicines, change doses, or stop long-term treatment without advice.

\section*{Complications}
Complications depend on the underlying cause and may include progression, effects on other organs, treatment adverse effects, or reduced quality of life.

\section*{Prognosis and Self-Care}
Outcome depends on the cause, severity, complications, and response to treatment. Follow the care plan, keep appointments, avoid known triggers, and seek help if symptoms worsen.

\section*{Expanded clinical context}
Over the counter medicines is a medicine or medication-use topic. Interpretation should account for age, baseline health, medicines, comorbidities, goals, and access to care. This chapter supports discussion with a qualified clinician and does not establish a diagnosis.

\section*{Assessment and decision points}
Review the product, indication, dose, route, interactions, allergies, kidney or liver disease, pregnancy, other medicines, and current symptoms. Document the main concern, time course, functional impact, relevant risks, and red flags. Use targeted investigations when their result is likely to change management.

\section*{Management and follow-up}
Management may include dose adjustment, substitution, monitoring, adherence support, adverse-effect treatment, or supervised tapering. Do not share medicines. Explain expected improvement, when reassessment is needed, and how follow-up can be accessed. Avoid starting, stopping, or sharing prescription medicines without professional advice.

\section*{Safety-netting}
Seek urgent help for breathing difficulty, facial swelling, collapse, seizure, chest pain, suspected overdose, or serious allergy.

\section*{Practical prevention}
Keep one current medicine list, follow labels, check interactions with a pharmacist, and store medicines securely. Keep written instructions and seek review if the topic recurs, changes, or does not improve as expected.

\section*{Limitations}
This educational expansion should be checked against current local guidance and authoritative topic-specific sources.
\clearpage
